\documentclass[11pt,a4paper]{article}
\usepackage[margin=23mm]{geometry}
\usepackage{amsmath,amssymb}
\usepackage[hidelinks]{hyperref}
\setlength{\parindent}{0pt}
\setlength{\parskip}{6pt}
\raggedright
\begin{document}
{\Large BP-R Stage 3: approved priority and static contract}\\[4pt]
6 October 2026 --- source/static review; S3-G3 disposition pending

\textbf{Status.} Reference route A, upstream admission M.A and fixed junction
priority N.A are approved. The unavailable third source is excluded from exact
replication. The equations below specify interface accounting and local bounds;
they do not certify a completed reference solver, a transient compact closure,
clearance, control benefit or a gradient.

\section*{1. The approved receiving allocation}
All quantities here are rates in vehicle equivalents per second.
$D_r\ge0$ is ramp sending demand at the merge; $D_m\ge0$ is the sum of compatible
non-ramp incoming demands to the same shoulder-lane receiver; $R\ge0$ is the
available receiving rate consistent with the traffic field and existing moving
constraints. $p\in(0,1)$ is a fixed physical ramp priority, common across models
and control schedules. It is not the upstream metering rate.

The selected conservative node operator is
\begin{equation}
(f_r,f_m)=
\begin{cases}
(D_r,D_m),&D_r+D_m\le R,\\
\bigl(\min\{D_r,\max\{pR,R-D_m\}\},\ R-f_r\bigr),&D_r+D_m>R.
\end{cases}
\end{equation}
In the second branch the displayed first coordinate defines $f_r$ before the
second coordinate is evaluated. This follows the fixed-priority merge family
in Daganzo's 1994 network CTM original, Section 3.2. Grouping non-ramp inflows
and joining moving constraints are declared task adaptations.

Under scarcity, $D_r>R-D_m$ and $pR\in[0,R]$. Hence
\begin{equation}
\max\{0,R-D_m\}\le f_r\le\min\{D_r,R\},\qquad
0\le f_m=R-f_r\le D_m.
\end{equation}
Together with the unconstrained branch,
\begin{equation}
f_r+f_m=\min\{D_r+D_m,R\},\qquad 0\le f_b\le D_b\quad(b=r,m).
\end{equation}
Thus receiving is spent once and unused shares are reallocated. If both streams
can use their contested shares, the allocation is $pR$ and $(1-p)R$.

If compatible non-ramp movements have already budgeted requests $d_j$,
$D_m=\sum_jd_j$, their subdivision is
\begin{equation}
f_j=\frac{d_j}{D_m}f_m\quad(D_m>0);\qquad f_j=0\quad(D_m=0).
\end{equation}
This preserves request bounds and sums to $f_m$. Donor budgets must be formed
before subdivision; independent full-capacity requests are not admissible.
Other receiving lanes keep their source local allocation.

\newpage
\section*{2. Upstream admission and population ownership}
$W^k$ is waiting mass upstream of the finite ramp approach, $A^k$ is offered
arrival mass in a time bin, $r^k$ is the common permission rate, $R_a^k$ is
approach receiving and $\Delta t$ is the bin duration. A reviewable discrete
boundary transcription is
\begin{equation}
B^k=\min\{W^k+A^k,\ r^k\Delta t,\ R_a^k\Delta t\},\qquad
W^{k+1}=W^k+A^k-B^k.
\end{equation}
Admission $B^k$ is credited once to the ramp. It is not actual merge entry.
Stored approach traffic can enter the freeway faster than the contemporaneous
upstream command. All models must retain that propagation and count $W$.

For road/ramp control volumes $c$, let $n_c^k$ be continuous mass. Each internal
transfer $F_{cb}^k$ is debited from $c$ and credited to $b$. Let $I^k$ include
all offered external arrival mass, and $E^k$ include all sink exit mass. For
total waiting plus road/ramp inventory $M^k$,
\begin{equation}
M^k=W_r^k+W_m^k+\sum_cn_c^k,\qquad
M^{k+1}-M^k=I^k-E^k.
\end{equation}
Here $W_m$ includes any blocked mainline entrance population. Every internal
debit/credit cancels. Conservation alone does not establish receiving/storage
correctness; those also require the coupled flux law.

The supplied journal's lateral recipe sends mass from $(i,\ell)$ to $(i+1,b)$.
Counting a separate same-cell lateral transfer followed by another longitudinal
transfer for that mass would change the source update. Descriptor insertion
does not add mass to $M$.

For a particular actual entry/transfer stream, let $H(t)$ be its accumulated
continuous mass plus a declared initial event progress. The deterministic
source descriptor count can be expressed as
\begin{equation}
K(t)=\lfloor H(t)\rfloor-\lfloor H(0)\rfloor.
\end{equation}
This selects unit-threshold descriptors; it does not round $H$, the flow or
$n_c$. Event progress is public initial data, not an optimized phase. Threshold
crossing times and coincident events remain part of the reference transcription.

The 2004 antecedent, Section 3, gives the sending-lane speed as birth speed.
The final 2006 journal, Appendix A, then constrains realized motion using
\begin{equation}
v(t)=\min\{v_{\rm des}(t),v_{\rm dwn}(t)\},\qquad
a(v,x)=a_0(1-v/v_{\max})-gG(x).
\end{equation}
$v_{\rm dwn}$ is the immediately downstream KW-stream speed, not the lane-choice
look-ahead speed. Positions move with realized speed; desired speed advances
with the source acceleration law. $a_0,v_{\max},G$ are physical parameters, not
functions assigned arbitrarily from permission $r$.

\newpage
\section*{3. Complete system time and the remaining boundary}
Let $A(t)$ and $E(t)$ be cumulative offered arrivals and exits from the declared
start. For initial total population $M_0$,
\begin{equation}
M(t)=M_0+A(t)-E(t),\qquad
J_{\rm TST}=\int_0^{T_{\rm clear}}M(t)\,dt.
\end{equation}
$J_{\rm TST}$ has units vehicle-equivalent seconds. Mean remaining system time
is $J_{\rm TST}/P$ for common finite population $P$; pre-start elapsed time is
a common constant. An asymptotic convergent inventory integral with a certified
tail bound may replace finite $T_{\rm clear}$. No extra merge-work cost is added.

For schedules $a,b$ with common population, arrivals and initial inventory,
extend each cleared exit history constantly to a common terminal time $T$:
\begin{equation}
J_{\rm TST}^{a}-J_{\rm TST}^{b}
=-\int_0^T\bigl(E^a(t)-E^b(t)\bigr)\,dt.
\end{equation}
Identical full exits therefore mean identical TST even if queues are relocated.
Different mean capacity or one terminal count is not a complete-objective
comparison. This identity is bookkeeping, not a feasible-control witness.

\textbf{Clearance is not proved.} Finite offered arrivals and a positive draining
permission do not alone establish clearance of the adapted reference. The joined
law must justify eventual discharge: reachable exits, available receiving and
constraints that release or depart. A certified residual tail is an alternative;
dropping remaining mass or assuming an empty road is not.

\section*{4. What the source evidence does and does not finish}
The supplied journal and 2004 antecedent define traffic/event ingredients.
The 2003 moving-boundary report gives a cumulative-count coupling route; its
node/event integration here is not yet verified. A guessed acceleration-capacity
penalty cannot replace that route.

Existing compact alternatives are serious. The 2011 endogenous merge work and
2016 moment manuscript predict mean capacity under their assumptions, not a
matched transient state update for this upstream meter. The 2021 hysteretic CTM
uses a different receiving/control construction. Importing any of these without
qualification would confound representation with a changed task law.

An independent rival must initialize from its declared public projection and
evolve its own speed/mode/moment state. It cannot borrow the reference's future
entries or constraint history. Unknown accuracy of a bound law can be tested
later; a missing law cannot be silently delegated to an experiment.

\textbf{Result: partial/evidence-limited formulation.} Local bounds and accounting
are explicit. A complete joined reference/clearance and a strong independent
transient compact closure remain unqualified. This is neither a necessity proof
for finer state nor a sufficiency/impossibility theorem for compact models.
Stage 3 disposition is pending S3-G3; Stage 4 remains unopened.

Source versions and links are recorded in
\texttt{post\_priority\_evidence\_audit.md}; the plain-language interpretation
and remaining interfaces are in
\texttt{joined\_contract\_and\_reduction\_audit.md}.
\end{document}
